\documentclass[12pt,a4paper]{article}

\usepackage[utf8]{inputenc}
\usepackage[T1]{fontenc}
\usepackage[brazil]{babel}
\usepackage{amsmath}
\usepackage{enumitem}
\usepackage{geometry}
\usepackage{hyperref}

\geometry{margin=2.5cm}

\title{Protocolo de Implementação do Simulador do Paradoxo de Braess}
\author{Diego Trigo Araujo}
\date{\today}

\begin{document}

\maketitle

\section{Objetivo}

Este documento define a ordem de implementação, os requisitos técnicos e os
critérios de validação do código utilizado na pesquisa sobre o paradoxo de
Braess.

O objetivo principal do software é:

\begin{quote}
Modelar uma sub-rede viária como um multidígrafo, calcular uma aproximação do
equilíbrio de tráfego, remover conexões da rede e verificar se alguma remoção
reduz o tempo total de viagem do sistema.
\end{quote}

O desenvolvimento deverá priorizar:

\begin{itemize}
    \item correção matemática;
    \item preservação da estrutura de multidígrafo;
    \item separação entre etapas;
    \item testes automatizados;
    \item reprodutibilidade dos experimentos;
    \item rastreabilidade dos parâmetros utilizados.
\end{itemize}

\section{Escopo do software}

O software deverá implementar:

\begin{enumerate}
    \item obtenção ou carregamento de uma rede viária;
    \item preservação da rede como \texttt{MultiDiGraph};
    \item preparação dos atributos das arestas;
    \item representação de demanda origem-destino;
    \item funções de custo por aresta;
    \item cálculo de menor caminho;
    \item atribuição de tráfego;
    \item aproximação do equilíbrio de Wardrop;
    \item remoção controlada de conexões;
    \item comparação entre cenários;
    \item exportação de resultados.
\end{enumerate}

Não fazem parte do escopo inicial:

\begin{itemize}
    \item simulação dinâmica de filas;
    \item tráfego em tempo real;
    \item interface gráfica;
    \item modelagem microscópica de veículos;
    \item controle semafórico detalhado;
    \item previsão de trânsito.
\end{itemize}

\section{Estrutura planejada}

A estrutura inicial do projeto deverá seguir:

\begin{verbatim}
src/
  braess/
    graph/
    demand/
    costs/
    routing/
    equilibrium/
    experiments/
    metrics/
    io/

tests/
  unit/
  integration/
  validation/

experiments/
outputs/
scripts/
\end{verbatim}

Cada módulo deverá possuir responsabilidade única.

\section{Regras gerais de implementação}

\begin{enumerate}
    \item A rede não deverá ser convertida para \texttt{DiGraph} durante o
    experimento principal.
    
    \item Cada aresta deverá ser identificada pela tripla
    \((u,v,k)\).
    
    \item Parâmetros experimentais não deverão ficar fixos dentro do código.
    
    \item Demanda, tolerância, função de custo e área analisada deverão ser
    definidos em arquivos de configuração.
    
    \item Cada etapa deverá possuir testes independentes.
    
    \item Um cenário não convergente não poderá ser classificado como resultado
    válido.
    
    \item Toda execução deverá salvar parâmetros, métricas e versão do código.
\end{enumerate}

\section{Fase 1 --- Preparação do experimento}

\subsection{Tarefa 1.1 --- Criar arquivo de configuração}

Criar um arquivo YAML para cada experimento.

Exemplo:

\begin{verbatim}
experiments/
  synthetic_braess.yaml
  sao_jose.yaml
\end{verbatim}

O arquivo deverá conter:

\begin{itemize}
    \item identificação do experimento;
    \item origem dos dados;
    \item coordenadas;
    \item raio;
    \item função de custo;
    \item algoritmo de equilíbrio;
    \item tolerância;
    \item limite de iterações;
    \item matriz origem-destino;
    \item unidade de remoção;
    \item diretório de saída.
\end{itemize}

\textbf{Critério de aceite:}

Nenhum parâmetro metodológico relevante deverá depender de uma constante
espalhada no código.

\subsection{Tarefa 1.2 --- Definir identidade de aresta}

Criar uma estrutura imutável para representar arestas:

\begin{verbatim}
@dataclass(frozen=True, slots=True)
class EdgeId:
    u: int
    v: int
    key: int
\end{verbatim}

\textbf{Critério de aceite:}

Duas arestas paralelas entre os mesmos nós deverão poder receber fluxos e
custos diferentes.

\subsection{Tarefa 1.3 --- Definir formato de resultados}

Definir previamente os arquivos gerados por cada execução:

\begin{itemize}
    \item \texttt{metadata.json};
    \item \texttt{iterations.csv};
    \item \texttt{edge\_flows.csv};
    \item \texttt{summary.csv};
    \item \texttt{removed\_connections.csv}.
\end{itemize}

\textbf{Critério de aceite:}

Uma execução deverá poder ser analisada posteriormente sem necessidade de
executar novamente o código.

\section{Fase 2 --- Rede e pré-processamento}

\subsection{Tarefa 2.1 --- Carregar a rede como multidígrafo}

O módulo deverá:

\begin{itemize}
    \item carregar dados do OpenStreetMap;
    \item utilizar \texttt{networkx.MultiDiGraph};
    \item preservar direção e arestas paralelas;
    \item salvar o grafo bruto antes do pré-processamento.
\end{itemize}

\textbf{Critério de aceite:}

O tipo do objeto final deverá permanecer \texttt{MultiDiGraph}.

\subsection{Tarefa 2.2 --- Validar topologia}

A validação deverá verificar:

\begin{itemize}
    \item número de nós;
    \item número de arestas;
    \item componentes fortemente conectadas;
    \item arestas paralelas;
    \item laços;
    \item atributos ausentes;
    \item presença de caminhos entre pares origem-destino.
\end{itemize}

\textbf{Critério de aceite:}

Deverá ser gerado um relatório contendo os problemas encontrados.

\subsection{Tarefa 2.3 --- Preparar velocidades}

Para cada aresta, determinar:

\begin{itemize}
    \item velocidade obtida diretamente da fonte;
    \item velocidade estimada por classe viária;
    \item origem do valor utilizado.
\end{itemize}

O atributo final deverá ser:

\begin{verbatim}
speed_kph
speed_source
\end{verbatim}

\textbf{Critério de aceite:}

Nenhuma aresta válida deverá permanecer sem velocidade.

\subsection{Tarefa 2.4 --- Calcular tempo de fluxo livre}

O tempo deverá ser calculado por:

\begin{equation}
    t_e^0 = \frac{l_e}{v_e}
\end{equation}

O valor deverá ser armazenado no atributo:

\begin{verbatim}
free_flow_time
\end{verbatim}

\textbf{Critério de aceite:}

Todo tempo de fluxo livre deverá ser positivo.

\subsection{Tarefa 2.5 --- Estimar capacidade}

A capacidade deverá considerar:

\begin{itemize}
    \item tipo de via;
    \item quantidade de faixas;
    \item sentido;
    \item valor padrão quando houver dado ausente;
    \item origem da regra aplicada.
\end{itemize}

Os atributos deverão ser:

\begin{verbatim}
capacity
capacity_source
\end{verbatim}

\textbf{Critério de aceite:}

Toda aresta utilizada na simulação deverá possuir capacidade positiva.

\section{Fase 3 --- Demanda}

\subsection{Tarefa 3.1 --- Criar modelo origem-destino}

Criar uma estrutura:

\begin{verbatim}
@dataclass(frozen=True, slots=True)
class ODPair:
    origin: int
    destination: int
    demand: float
\end{verbatim}

\textbf{Critério de aceite:}

A demanda deverá ser positiva e os nós deverão existir no grafo.

\subsection{Tarefa 3.2 --- Validar matriz origem-destino}

Verificar:

\begin{itemize}
    \item existência dos nós;
    \item conectividade;
    \item demanda total;
    \item duplicidade de pares;
    \item unidade utilizada.
\end{itemize}

\textbf{Critério de aceite:}

Nenhum par desconectado deverá entrar no solver sem ser explicitamente
sinalizado.

\section{Fase 4 --- Funções de custo}

\subsection{Tarefa 4.1 --- Criar interface comum}

A interface deverá expor:

\begin{verbatim}
travel_time(...)
integral(...)
\end{verbatim}

A integral será utilizada pelo algoritmo de Frank--Wolfe.

\subsection{Tarefa 4.2 --- Implementar função linear}

Implementar:

\begin{equation}
    t_e(x_e) = \alpha_e + \beta_e x_e
\end{equation}

\textbf{Critério de aceite:}

A função deverá reproduzir exatamente os valores do exemplo sintético.

\subsection{Tarefa 4.3 --- Implementar função BPR}

Implementar:

\begin{equation}
    t_e(x_e)
    =
    t_e^0
    \left[
        1 +
        \alpha
        \left(
            \frac{x_e}{c_e}
        \right)^\beta
    \right]
\end{equation}

\textbf{Critérios de aceite:}

\begin{itemize}
    \item \(t_e(0)=t_e^0\);
    \item o tempo não diminui quando o fluxo aumenta;
    \item a capacidade deve ser positiva;
    \item a integral deve possuir teste automatizado.
\end{itemize}

\section{Fase 5 --- Roteamento}

\subsection{Tarefa 5.1 --- Menor caminho no multidígrafo}

Implementar uma função que retorne:

\begin{verbatim}
list[EdgeId]
\end{verbatim}

e não apenas uma lista de nós.

\textbf{Critério de aceite:}

Em um grafo com duas arestas paralelas, a aresta com menor custo deverá ser
selecionada corretamente pela chave \texttt{key}.

\subsection{Tarefa 5.2 --- Atribuição All-or-Nothing}

Para cada par origem-destino:

\begin{enumerate}
    \item calcular o menor caminho;
    \item atribuir toda a demanda ao caminho;
    \item acumular o fluxo nas arestas;
    \item retornar o vetor auxiliar de fluxos.
\end{enumerate}

\textbf{Critério de aceite:}

A demanda total atribuída deverá ser igual à demanda total de entrada.

\section{Fase 6 --- Equilíbrio}

\subsection{Tarefa 6.1 --- Implementar métricas}

Implementar:

\begin{itemize}
    \item tempo total do sistema;
    \item tempo médio;
    \item objetivo de Beckmann;
    \item lacuna relativa;
    \item número de iterações.
\end{itemize}

\subsection{Tarefa 6.2 --- Implementar MSA}

O MSA será utilizado como solver inicial e referência.

\textbf{Critério de aceite:}

O algoritmo deverá convergir no exemplo sintético dentro de tolerância
definida.

\subsection{Tarefa 6.3 --- Implementar Frank--Wolfe}

Cada iteração deverá executar:

\begin{enumerate}
    \item atualização dos custos;
    \item atribuição All-or-Nothing;
    \item cálculo da direção;
    \item busca do tamanho do passo;
    \item atualização do fluxo;
    \item cálculo da convergência.
\end{enumerate}

A atualização será:

\begin{equation}
    \mathbf{x}^{(n+1)}
    =
    \mathbf{x}^{(n)}
    +
    \lambda_n
    \left(
        \mathbf{y}^{(n)} - \mathbf{x}^{(n)}
    \right)
\end{equation}

\textbf{Critério de aceite:}

O solver deverá retornar:

\begin{itemize}
    \item fluxo final;
    \item histórico de iterações;
    \item estado de convergência;
    \item lacuna relativa final;
    \item tempo total do sistema.
\end{itemize}

\section{Fase 7 --- Validação sintética}

\subsection{Tarefa 7.1 --- Criar rede clássica de Braess}

Criar um grafo com os nós:

\begin{verbatim}
O, A, B, D
\end{verbatim}

e as arestas definidas no trabalho.

\subsection{Tarefa 7.2 --- Validar cenário sem conexão}

Com demanda de \(4000\), o resultado esperado é aproximadamente:

\begin{equation}
    65 \text{ minutos}
\end{equation}

\subsection{Tarefa 7.3 --- Validar cenário com conexão}

Com demanda de \(4000\), o resultado esperado é aproximadamente:

\begin{equation}
    80 \text{ minutos}
\end{equation}

\textbf{Critério de aceite geral:}

Os dois cenários deverão ser reproduzidos por testes automatizados dentro de
uma tolerância explicitamente definida.

\section{Fase 8 --- Remoção de conexões}

\subsection{Tarefa 8.1 --- Definir unidade de remoção}

A implementação deverá diferenciar:

\begin{itemize}
    \item aresta direcionada;
    \item segmento bidirecional;
    \item rua física composta por múltiplas arestas.
\end{itemize}

\subsection{Tarefa 8.2 --- Criar cenário-base}

O cenário-base deverá salvar:

\begin{itemize}
    \item fluxos;
    \item custos;
    \item tempo total;
    \item tempo médio;
    \item convergência;
    \item parâmetros.
\end{itemize}

\subsection{Tarefa 8.3 --- Executar remoções}

Para cada conexão:

\begin{enumerate}
    \item copiar o grafo original;
    \item remover a conexão;
    \item validar conectividade;
    \item recalcular o equilíbrio;
    \item comparar com o cenário-base;
    \item salvar o resultado.
\end{enumerate}

\section{Fase 9 --- Detecção}

Uma remoção será classificada como candidata quando:

\begin{enumerate}
    \item o cenário-base convergir;
    \item o cenário modificado convergir;
    \item a rede permanecer válida;
    \item a demanda permanecer igual;
    \item o tempo total diminuir;
    \item a redução superar a tolerância numérica.
\end{enumerate}

A diferença será:

\begin{equation}
    \Delta T
    =
    T_{\mathrm{base}}
    -
    T_{\mathrm{rem}}
\end{equation}

e a melhoria relativa:

\begin{equation}
    I
    =
    \frac{
        T_{\mathrm{base}}
        -
        T_{\mathrm{rem}}
    }{
        T_{\mathrm{base}}
    }
\end{equation}

\section{Fase 10 --- Análise de sensibilidade}

Executar variações em:

\begin{itemize}
    \item demanda;
    \item capacidade;
    \item parâmetros da função de custo;
    \item tolerância;
    \item pares origem-destino;
    \item área da rede.
\end{itemize}

\textbf{Critério de aceite:}

Uma conexão deverá ser considerada mais robusta quando a melhoria aparecer
em múltiplos cenários e não apenas em uma configuração isolada.

\section{Ordem de implementação}

A ordem recomendada é:

\begin{enumerate}
    \item configuração;
    \item \texttt{EdgeId};
    \item função linear;
    \item menor caminho em multidígrafo;
    \item All-or-Nothing;
    \item métricas;
    \item MSA;
    \item validação sintética;
    \item Frank--Wolfe;
    \item função BPR;
    \item preparação da rede real;
    \item capacidade;
    \item demanda urbana;
    \item cenário-base;
    \item remoção de conexões;
    \item análise de sensibilidade.
\end{enumerate}

\section{Definição de pronto}

O projeto será considerado funcional quando:

\begin{enumerate}
    \item reproduzir o exemplo sintético;
    \item preservar corretamente arestas paralelas;
    \item conservar os fluxos;
    \item demonstrar convergência;
    \item executar um cenário-base urbano;
    \item remover conexões de forma controlada;
    \item gerar resultados reproduzíveis;
    \item exportar métricas suficientes para análise.
\end{enumerate}

\end{document}